\documentclass[11pt]{article}
\usepackage[margin=1in]{geometry}
\usepackage{amsmath,amssymb,amsthm,mathtools}
\usepackage{booktabs}
\usepackage{enumitem}
\usepackage{hyperref}

\newcommand{\WARP}{\textsc{Warp}}
\newcommand{\Hist}{\mathsf{Hist}}
\newcommand{\Tr}{\mathbf{Tr}}
\newcommand{\WState}{\mathsf{WState}}
\newcommand{\Apply}{\mathsf{Apply}}
\newcommand{\Dist}{\mathsf{Dist}}
\newcommand{\DL}{\mathsf{DL}}

\theoremstyle{definition}
\newtheorem{observation}{Observation}
\newtheorem{principle}{Principle}
\newtheorem{speculation}{Speculation}

\title{Observer Geometry, Profunctor Optics, and the Effect Architecture\\[6pt]
\large Design Notes from the git-warp Effect Emission Session\\[3pt]
\normalsize 2026-03-30}
\author{}
\date{}

\begin{document}
\maketitle

\section{Context}

These notes record the design discoveries from a working session on
outbound effect emission in \texttt{git-warp}. The session began as an
implementation task (``add substrate support for effect emission and
delivery observation'') and evolved into a clarification of the
observer/participant boundary, the role of profunctor optics in
\WARP{} rewriting, and the architectural implications for the
Continuum dual-runtime model.

The primary sources are:
\begin{itemize}[leftmargin=2em]
\item WARP Papers I--IV (AI\textOmega{}N Foundations Series)
\item Observer Geometry I and IV (OG Series)
\item WARP Optics: A Working Mapping (Paper~07, working draft)
\item \texttt{xyph/design/effect-emission-alignment.md}
\item \texttt{xyph/design/substrate-alignment.md}
\end{itemize}

\section{The Three-Layer Invariant}

The central discovery is a strict three-layer separation that must be
preserved for deterministic replay to hold.

\begin{principle}[Three-Layer Separation]
\label{princ:three-layer}
The system has three layers with distinct causal status:
\begin{enumerate}[leftmargin=2em]
\item \textbf{Graph layer} (causal, deterministic). Inert structure:
  nodes, edges, properties. Written by participants (writers). Given
  the same initial state and patches, materialization always produces
  the same result. The graph never acts.

\item \textbf{Observer layer} (non-causal, deterministic for fixed
  input). Read-only projections over worldlines. An observer
  $O : \Hist(\mathcal{U}, R) \to \Tr$ maps derivation paths to trace
  values (Paper~IV, Definition~3.3). Observers accumulate state and
  emit structural descriptions (OG-I, Definition~3), but never write
  to the causal graph.

\item \textbf{Application layer} (non-causal, policy-governed).
  Consumes observer descriptions. Decides external actions: call a
  webhook, log to console, suppress during replay. Bridges observer
  outputs to writers when recording is desired.
\end{enumerate}
\end{principle}

\begin{observation}[Observers Must Not Write]
\label{obs:no-write}
If an observer writes to the causal graph, the graph state becomes
dependent on \emph{who is watching}. Two replicas with different
observers would produce different graph states. Deterministic replay
is broken. This is the observer-participant boundary: observers
describe, participants act. Crossing this boundary destroys the
invariant that makes \WARP{} replay sound.
\end{observation}

\begin{observation}[The Graph Is Inert]
\label{obs:inert}
Nothing about graph structure causes anything to happen. Effect nodes
exist in the graph because a writer put them there. Delivery edges
exist because a writer recorded them. The graph is the \emph{record}
of what happened, not the \emph{cause} of what happens. Causation
flows through application logic that bridges observer outputs to
writer actions.
\end{observation}

\section{The OG-I Structural Observer}

OG-I defines the structural observer as a five-tuple:
\[
S = (O, \mathcal{B}_S, M_S, K_S, E_S)
\]
where:
\begin{itemize}[leftmargin=2em]
\item $O : \Hist(\mathcal{U}, R) \to \Tr_S$ is the raw projection
  (what is seen),
\item $\mathcal{B}_S$ is the observer basis (what distinctions are
  natively expressible),
\item $M_S$ is the observational state space (memory),
\item $K_S : M_S \times \Tr_S \to M_S$ is the update rule (how
  observations accumulate),
\item $E_S : M_S \to \widehat{\Tr}_S$ is the emission map (what
  structural description is externally available after accumulation).
\end{itemize}

\begin{observation}[Emission Is Description, Not Side Effect]
The emission map $E_S$ produces \emph{accumulated structural
descriptions}---not external side effects. It is the observer's
output in $\widehat{\Tr}_S$ (accumulated trace space), which is
external to the graph. Console output, webhook calls, and debug
displays are application-level actions taken \emph{in response to}
observer emissions, not part of the emission itself.
\end{observation}

\section{Effect Architecture}

Given the three-layer invariant, the effect delivery architecture is:

\begin{enumerate}[leftmargin=2em]
\item A \textbf{writer} (participant) writes effect nodes
  (\texttt{effect:*}) to the graph as part of a patch. This is a
  causal act.
\item An \textbf{observer} watches the worldline through a projection
  that matches effect nodes. It accumulates state
  (``these effects are new'') and emits structural descriptions.
\item \textbf{Application logic} receives the observer's emission and
  decides what to do: deliver externally, suppress (replay mode),
  record delivery via a writer.
\item If the application wants delivery recorded in the graph, a
  \textbf{writer} writes a subsequent patch. The observer does not
  write---the application bridges the observer output to the writer.
\end{enumerate}

During replay:
\begin{itemize}[leftmargin=2em]
\item The graph is the same (deterministic materialization).
\item The observer sees the same effect nodes (same projection, same
  input).
\item Application policy suppresses external delivery.
\item Suppression is visible in the application's delivery log
  (observer emission trace), not in the graph.
\end{itemize}

\section{The Lens Question}

In the current \texttt{git-warp} codebase, \texttt{Lens} is a simple
projection configuration:
\[
\text{Lens} = \{\, \text{match},\; \text{expose},\; \text{redact} \,\}
\]
This is the parameterization of the observer's raw projection $O$.
It is not a separate first-class entity---it is observer configuration.

The WARP Optics working paper (Paper~07) proposes a richer notion of
\emph{Lens} drawn from profunctor optics: a bidirectional
$(get, set)$ pair that focuses on part of a whole, transforms the
focused part, and reassembles the whole. In this reading:

\begin{itemize}[leftmargin=2em]
\item The \textbf{get} side is the observer's projection (read-only,
  deterministic).
\item The \textbf{set} side is the writer's rewrite (causal,
  participatory, changes the graph).
\item The lens is the \emph{structural declaration} connecting them:
  the focus shape. Different actors use different sides.
\end{itemize}

\begin{principle}[Lens As Declaration, Not Actor]
A lens is not an actor. It is a contract that declares what part of
the whole is being focused on. The observer reads through it. The
writer writes through it. The lens does not read or write---it
\emph{is} the shape of the focus.
\end{principle}

For the current implementation, the simple projection configuration
(\texttt{match/expose/redact}) is sufficient. The full optic lens
becomes relevant when Wesley compiles rewrite shapes that need
multiple interpretations.

\section{Profunctor Optics and \WARP{} Rewrites}

The profunctor representation of an optic is:
\[
\mathcal{O} : p\,a\,b \to p\,s\,t
\]
where $s$ is the whole before transformation, $t$ is the whole after,
$a$ is the focused part before, $b$ is the focused part after, and
$p$ is an abstract transformation context.

The power of this representation is that $p$ is not hardcoded. The
\emph{same} declared rewrite shape, instantiated with different
profunctors, yields different interpretations:

\begin{center}
\begin{tabular}{@{}ll@{}}
\toprule
Profunctor instantiation $p$ & Interpretation \\
\midrule
$p$ = execute & Apply the rewrite (Echo runtime) \\
$p$ = merge & CRDT-merge the patch (git-warp runtime) \\
$p$ = explain & Produce a debugger trace (warp-ttd) \\
$p$ = footprint & Derive the read/write/delete sets \\
$p$ = conflict-cost & Estimate rulial distance between competing rewrites \\
$p$ = witness & Extract the minimal reversibility residue \\
$p$ = compile & Emit a Wesley-compiled artifact \\
\bottomrule
\end{tabular}
\end{center}

\begin{observation}[Compile Once, Interpret Many]
\label{obs:compile-once}
Wesley compiles a declared rewrite shape from a GraphQL-like schema
into typed code with footprint declarations. The profunctor
parametricity means the \emph{same compiled shape} can be
instantiated as Echo execution, git-warp merge, warp-ttd explanation,
or conflict analysis---without redeclaring the rewrite for each use
case.
\end{observation}

\subsection{Speculative Applications}

\begin{speculation}[Observer-Relative Conflict as Enriched Optic Composition]
Two optics $\mathcal{O}_1$ and $\mathcal{O}_2$ with overlapping focus
regions cannot compose freely---their footprints interfere. The
interference relation (Paper~II, Definition~4.7; OG-IV, Definition~7)
is the obstruction to composition. In the optic frame, this could be
expressed as a failure of the profunctor composition law, with the
conflict artifact being the explicit witness of the failure. Rulial
distance between two conflicting rewrites then becomes the cost of
translating the composed-with-conflict optic into a conflict-free one.
\end{speculation}

\begin{speculation}[Speculative Lanes as Optic Forks]
A strand in \texttt{git-warp} is a speculative copy-on-write lane. In
optic terms, creating a strand is forking the residual $M$: the
focused region is shared, but each strand carries its own hidden
context. Braiding strands is a controlled merge of residuals.
Transfer/collapse is selecting one residual as canonical and discarding
alternatives. The coend quotient (Paper~07, \S4) identifies equivalent
residual presentations---so two strands that took different paths to
the same focused state would be identified up to the relevant
quotient.
\end{speculation}

\begin{speculation}[Time-Travel as Optic Inversion]
If a \WARP{} optic $\Omega = (\pi, \phi, \rho, \omega, \sigma)$
carries a witness $\omega$ sufficient for local reversibility, then
time-travel (seek to a previous coordinate) is the repeated
application of $\rho^{-1}$ under witness $\omega$. The current
\texttt{TickReceipt} contains more than the minimal witness---it is
the operational envelope. Splitting it into witness (semantic minimum
for inversion) and receipt (operational metadata) would make time-travel
formally an optic inversion rather than a full rematerialization.
\end{speculation}

\begin{speculation}[Effect Delivery as Profunctor Instantiation]
If effect emission is modeled as an optic over effect nodes (focus on
\texttt{effect:*}, transform by adding delivery edges), then the
delivery policy is a profunctor instantiation:
\begin{itemize}[leftmargin=2em]
\item $p$ = deliver-live: call the external adapter, write
  \texttt{delivered\_to} edge with \texttt{outcome: delivered}.
\item $p$ = deliver-replay: skip the external call, write
  \texttt{delivered\_to} edge with \texttt{outcome: suppressed}.
\item $p$ = deliver-inspect: dry-run, produce explanation trace
  without writing.
\end{itemize}
Same focus shape (the effect node), different transformation behavior.
\end{speculation}

\begin{speculation}[Wesley-Generated Observer Projections]
If Wesley compiles entity types from a schema, it knows which
properties each entity type declares. An observer projection matching
\texttt{effect:*} could be automatically generated by Wesley to
include exactly the declared properties of the \texttt{Effect} entity
type, with type-safe access. The current string-based
\texttt{match/expose/redact} lens would be replaced by a
schema-derived, type-checked projection.
\end{speculation}

\begin{speculation}[Cross-Runtime Provenance via Shared Optics]
If Echo and git-warp share the same Wesley-compiled optic shapes, then
a rewrite executed in Echo and the ``same'' rewrite merged in git-warp
carry compatible provenance. The warp-ttd debugger can display a
unified timeline where some ticks were executed by Echo (real-time,
Rust) and others were merged by git-warp (eventual, TypeScript). The
observer sees a single worldline; the translator cost between the two
runtimes' traces is bounded by the codec overhead.
\end{speculation}

\section{Continuum: The Dual-Runtime Model}

The Continuum architecture emerges from unifying Echo and git-warp
through Wesley-compiled schemas and profunctor optics:

\begin{center}
\begin{tabular}{@{}lll@{}}
\toprule
Component & Runtime & Role \\
\midrule
Echo & Rust & High-performance deterministic real-time execution \\
git-warp & TypeScript/Git & Offline-first CRDT storage with
  replication \\
Wesley & Compiler & Schema $\to$ types + footprints + binary codecs \\
warp-ttd & Observer & Time-travel debugger across both runtimes \\
\bottomrule
\end{tabular}
\end{center}

The same \WARP{} graph state can be:
\begin{itemize}[leftmargin=2em]
\item loaded from git-warp into Echo for real-time simulation,
\item committed from Echo back to git-warp for durable replication,
\item debugged through warp-ttd regardless of which runtime executed
  each tick.
\end{itemize}

Binary compatibility is ensured by Wesley-generated codecs. Type
safety is ensured by Wesley-generated types. Footprint correctness
is ensured by Wesley-derived footprint declarations. The profunctor
optic is the mechanism: same declared shape, different runtime
interpretation.

\begin{observation}[Two Runtimes, One Observer Geometry]
The observer geometry (Papers~IV, OG-I, OG-IV) applies uniformly
across both runtimes. Rulial distance between an Echo trace and a
git-warp trace of the ``same'' computation is bounded by the
translator overhead of the Wesley codec---which is a constant
independent of computation length. The two runtimes are
\emph{observer-equivalent up to codec cost}.
\end{observation}

\section{What We Shipped}

The implementation session produced application-level effect delivery
tooling for \texttt{git-warp}:

\begin{itemize}[leftmargin=2em]
\item \textbf{Types}: \texttt{EffectEmission}, \texttt{DeliveryObservation},
  \texttt{DeliveryLens} (application-level delivery policy).
\item \textbf{Port}: \texttt{EffectSinkPort} (hexagonal adapter interface
  for external delivery).
\item \textbf{Services}: \texttt{EffectPipeline} (orchestration),
  \texttt{MultiplexSink} (fan-out composite).
\item \textbf{Adapters}: \texttt{ConsoleEffectSink},
  \texttt{ChunkEffectSink}, \texttt{NoOpEffectSink}.
\item \textbf{WarpCore integration}: \texttt{emit()},
  \texttt{effectEmissions}, \texttt{deliveryObservations},
  \texttt{deliveryLens} get/set, \texttt{open()} options.
\end{itemize}

This is \emph{host-domain infrastructure}---basic tooling any \WARP{}
application needs for managed outbound delivery. It is not substrate.
The substrate is the graph, the observer machinery, and the worldline
model, which already existed.

The key design insight: \texttt{warp-ttd} does not need this tooling.
\texttt{warp-ttd} is itself an observer. It creates its own observer
over the worldline and produces its own traces. The ``sync point''
between \texttt{warp-ttd} and \texttt{git-warp} is the existing
observer/worldline/query machinery, not the effect pipeline.

\section{Open Questions}

\begin{enumerate}[leftmargin=2em]
\item What are the first \WARP{} optic laws worth enforcing in code?
\item Can the current \texttt{TickReceipt} be split into a minimal
  witness and an operational shell?
\item What does the coend quotient concretely look like for \WARP{}
  rewrite residuals?
\item Can Wesley compile a declared rewrite shape and emit multiple
  profunctor interpretations automatically?
\item Should effect node schemas be declared in Wesley's schema
  language from the start, or defined as a convention and formalized
  later?
\item What is the codec cost bound for Echo $\leftrightarrow$
  git-warp translation, and does it match the theoretical prediction
  of observer-equivalence up to constant overhead?
\end{enumerate}

\section{One-Sentence Summary}

The graph is inert, observers describe, application logic acts, and
the profunctor optic is the mechanism that lets one declared rewrite
shape be executed by Echo, merged by git-warp, and explained by
warp-ttd---which is what makes Continuum a single system with two
runtimes rather than two systems with a bridge.

\end{document}
